\documentclass[uplatex,a4paper,11pt]{jsarticle}

\usepackage[dvipdfmx]{graphicx}
\usepackage{amsmath, amssymb}
\usepackage{geometry}
\usepackage{hyperref}
\usepackage{here}
\usepackage{bookmark}

\geometry{margin=25mm}
\renewcommand{\baselinestretch}{1.2}

\title{Homework 5a}
\author{学籍番号: 1270374 \\ 名前: 松本 奈生}
\date{\today}

\begin{document}
\maketitle

%% ============================================================
\section*{Question 1}
%% ============================================================

$A_{TM}$ が決定可能であると仮定する．
すなわち，任意の入力 $\langle M, w \rangle$ に対して，
$M$ が $w$ を受理するかどうかを必ず判定する Turing 機械 $H$ が存在するとする．

このとき，言語
\[
L_D = \{ \langle M \rangle \mid M は \langle M \rangle を受理しない \}
\]
を決定する機械 $D$ を構成する．

$D$ は入力 $\langle M \rangle$ に対して以下のように動作する：

\begin{enumerate}
    \item $H(\langle M, \langle M \rangle \rangle)$ を実行する
    \item $H$ が accept した場合，reject する
    \item $H$ が reject した場合，accept する
\end{enumerate}

$H$ は必ず停止するため，$D$ も必ず停止する．
したがって $D$ は $L_D$ を決定する．

しかし Homework 4a で示したように，$L_D$ は認識不可能であるため，
決定可能であることはあり得ない．

よって仮定は誤りであり，$A_{TM}$ は決定不能である．

---

%% ============================================================
\section*{Question 2}
%% ============================================================

$A_{TM}$ が Turing 認識可能であることを示す．

$A_{TM} = \{ \langle M, w \rangle \mid M が w を受理する \}$ に対して，
以下の Turing 機械 $R$ を構成する：

\begin{enumerate}
    \item 入力 $\langle M, w \rangle$ を受け取る
    \item $M$ を入力 $w$ 上でシミュレートする
    \item $M$ が accept した場合，accept する
    \item $M$ が reject した場合，reject する
\end{enumerate}

もし $M$ が $w$ 上で無限ループする場合，$R$ も停止しない．

したがって，
\begin{itemize}
    \item $w \in A_{TM}$ のとき，必ず accept する
    \item $w \notin A_{TM}$ のとき，reject またはループ
\end{itemize}

となるので，$A_{TM}$ は Turing 認識可能である．

---

%% ============================================================
\section*{Question 3}
%% ============================================================

\subsection*{1.}

\[
\left[
\frac{ab}{abab}
\right]
\left[
\frac{ab}{abab}
\right]
\left[
\frac{aba}{b}
\right]
\left[
\frac{b}{a}
\right]
\left[
\frac{b}{a}
\right]
\left[
\frac{aa}{a}
\right]
\left[
\frac{aa}{a}
\right]
\]

---

\subsection*{2.}

与えられた PCP：
\[
\left\{
\left[\frac{ab}{abab}\right],\ 
\left[\frac{b}{a}\right],\ 
\left[\frac{aba}{b}\right]
\right\}
\]

この問題に match が存在しないことを示す．

各ドミノについて，上段と下段の長さの差を考える：

\[
\begin{aligned}
\left[\frac{ab}{abab}\right] &: \quad 2 - 4 = -2 \\
\left[\frac{b}{a}\right] &: \quad 1 - 1 = 0 \\
\left[\frac{aba}{b}\right] &: \quad 3 - 1 = +2
\end{aligned}
\]

任意の列に対して，上段と下段が一致するためには，
全体の長さが一致する必要があるため，これらの差の総和は 0 でなければならない．

したがって，ドミノ $\left[\frac{ab}{abab}\right]$ と $\left[\frac{aba}{b}\right]$ は同じ回数だけ使用されなければならない．

このとき，ドミノ $\left[\frac{b}{a}\right]$ は長さの差に影響を与えないため，
任意の回数挿入できる．

しかし，文字列として一致するためには，
途中までの文字列が常に一致している必要がある

ここで，各ドミノの先頭文字に注目すると：

\begin{itemize}
    \item $\left[\frac{ab}{abab}\right]$ は，上段・下段ともに $\mathtt{a}$ で始まる
    \item $\left[\frac{aba}{b}\right]$ は，上段は $\mathtt{a}$，下段は $\mathtt{b}$ で始まる
    \item $\left[\frac{b}{a}\right]$ は，上段は $\mathtt{b}$，下段は $\mathtt{a}$ で始まる
\end{itemize}

どのドミノを先頭に選んでも，上段と下段の最初の文字が一致しない場合が生じるか，
または途中で不一致が発生する．

したがって，いかなる列をとっても，上段と下段の文字列が完全に一致することはない．

ゆえに，この PCP インスタンスには match は存在しない．

---

%% ============================================================
\section*{Question 4}
%% ============================================================

$ALL_{TM}$ が決定不能であることを示す．

\[
ALL_{TM} = \{ \langle M \rangle \mid L(M) = \Sigma^* \}
\]

$A_{TM}$ からの帰着により証明する．

---

$A_{TM}$ の入力 $\langle M, w \rangle$ に対して，
新しい Turing 機械 $M'$ を次のように構成する：

\begin{enumerate}
    \item 入力 $x$ を受け取る
    \item $M(w)$ をシミュレートする
    \item $M(w)$ が accept した場合，accept
    \item それ以外（reject またはループ）の場合，reject
\end{enumerate}

---

このとき：

\begin{itemize}
    \item $M$ が $w$ を受理するなら，$M'$ はすべての入力を受理する
    \item $M$ が $w$ を受理しないなら，$M'$ はどの入力も受理しない
\end{itemize}

したがって：

\[
\langle M, w \rangle \in A_{TM} \iff \langle M' \rangle \in ALL_{TM}
\]

---

ここで $ALL_{TM}$ が決定可能と仮定すると，
この変換により $A_{TM}$ も決定可能になってしまう．

しかし $A_{TM}$ は決定不能であるため矛盾．

---

よって，

\[
ALL_{TM} \text{ は undecidable である．}
\]

---

\end{document}
